\documentclass[10pt,a4paper]{amsart}
\usepackage[margin=19mm]{geometry}
\usepackage{amsmath,amssymb,mathtools}
\usepackage[scr=rsfs,cal=euler]{mathalfa}
\usepackage{xfrac}
\usepackage[unicode,hidelinks,bookmarks=false]{hyperref}
\newcommand{\ie}{i.e.\ }
\newcommand{\cf}{cf.\ }
\newcommand{\resp}{respectively\ }
\newcommand{\Subsection}{\subsection}
\providecommand{\nobreakdash}{\nobreak\hbox{-}}
\setlength{\emergencystretch}{2em}
\multlinegap0pt
\usepackage{amssymb}
\usepackage{bm}
\usepackage{xcolor}
\newtheorem{remark}{Remark}
\newtheorem{lemma}{Lemma}
\newtheorem{corollary}{Corollary}
\title{Estimates on elliptic equations that hold only where the Hessian is large}
\author{Amit Kumar Acharya, Ram Baran Verma}
\date{Source: arXiv:2607.17045v1 (CC BY 4.0)}
\begin{document}
\maketitle

\noindent\emph{Source and licence.} Estimates on elliptic equations that hold only where the Hessian is large. Version arXiv:2607.17045v1 (CC BY 4.0), distributed by arXiv under the Creative Commons Attribution 4.0 International licence, http://creativecommons.org/licenses/by/4.0/, as recorded in the arXiv licence metadata for this exact version and shown on the abstract page https://arxiv.org/abs/2607.17045v1. Full licence notice: \texttt{licenses/arxiv-2607.17045-CC-BY-4.0.txt}.\par\smallskip

\noindent\textit{Editorial scope.} Counted unit: the article's corollary slepsilon (source0.tex lines 535--547). The article's complete original proof of it is delivered with it (source0.tex lines 548--562, one \texttt{proof} environment of 15 lines). No mathematics is written, repaired or re-derived: the statement and the proof are the article's own lines, in the article's own notation, and no retained line is substituted or re-flowed.  Retained uncounted as the source-local context the proof needs: lines 173--176; lines 508--521.  Nothing was omitted from any retained span, so the package carries no omission record.  No retained line contains a citation, so the wrapper carries no bibliography and no bibliography entry.  No retained line contains a figure environment, an image-inclusion command or drawing code, so no image is delivered and the package declares no asset dependency.  Editorial formatting only, in addition: an AMS \texttt{amsart} preamble that restates the article's own declarations and macros used by the retained lines, namely the environments \texttt{remark}, \texttt{lemma}, \texttt{corollary} on separate counters, together with the standard packages those lines need.  The retained environments print in this document as Remark 1, Lemma 1, Corollary 1 in order of appearance, whereas the article prints the same items with the numbering of its own full document.  The complete native source of the article as distributed in the arXiv e-print for this exact version is bundled unchanged as \texttt{sources/205527/source0.tex}.
\begin{remark}\label{remsc}
As a consequence of above relation we find that if $\mathcal{L}_{\kappa}^{-}\big(D^{2}u,Du\big)\leq C_{0}$~in~$\Omega$ then 
$\mathcal{L}_{\kappa}^{-}\big(D^{2}u_{r},Du_{r}\big)\leq \gamma C_{0}r^{2}$~in~$x_{0}+r\Omega$ \text{for} $r,\gamma\in[0,1].$ 
\end{remark}

\begin{lemma}\label{lepaa}
There exist two small $\kappa_{0},$ $\epsilon>0$ such that if $\kappa\leq \kappa_{0},$ then for any lower semi continuous function $u: B_{2}\longrightarrow\mathbb{R}$ such that
\begin{equation}\label{lepsilon}
\left\{
\begin{aligned}
u&\geq0~~\text{in}~~B_{2},\\
\mathcal{L}_{\kappa}^{-}\big(D^{2}u,Du\big)&\leq 1~~\text{in}~~B_{2},\\ 
\inf_{B_{1}}u&\leq 1,\\
\end{aligned}
\right.
s\end{equation}
we have 
\[|\{u>t\}\cap B_{1}|\leq \tilde{C}t^{-\epsilon}~~\text{for~all}~~t>0.\]
\end{lemma}

\begin{corollary}\label{slepsilon}
There exist positive constants $\tilde{k}_{0}>0,~\epsilon_{1}$ and such that if $k\leq \tilde{k}_{0},$ then for any $r\leq 1,$ $\alpha\in(0,1)$ and any lower semicontinuous function $u:~\overline{B_{2r}}\longrightarrow\mathbb{R}$ satisfying 
\begin{equation*}
\left\{
\begin{aligned}
u&\geq0~~~\text{in}~~B_{2r},\\
\mathcal{L}_{\kappa}^{-}\big(D^{2}u,Du\big)&\leq \epsilon_{1}~~~\text{in}~~B_{2r},\\ 
|\{u>r^{\alpha}\}\cap B_{r}|&\geq \frac{1}{2}|B_{r}|,\\
\end{aligned}
\right.
\end{equation*}
we have $u\geq \epsilon_{1}r^{\alpha}$ in $B_{r}.$
\end{corollary}

\begin{proof}
Let $\tau$ be a universal constant such that $\tilde{C}\tau^{-\epsilon}<|B_{1}|/2,$ where $\tilde{C}$ and $\epsilon>0$ are constant from Theorem \ref{lepsilon}. Now consider the function $v(x)=\tau r^{-\alpha}u(rx).$ This function satisfies
\begin{equation*}
\left\{
\begin{aligned}
v&\geq0~~~\text{in}~~B_{2},\\
\mathcal{L}_{\kappa\tau r^{2-\alpha}}^{-}\big(D^{2}v,Dv\big)&\leq\epsilon_{1}\tau r^{2-\alpha}~~~\text{in}~~B_{2},\\ 
|\{v>\tau\}\cap B_{1}|&\geq \frac{1}{2}|B_{1}|.\\
\end{aligned}
\right.
\end{equation*}
Now, we choose $\epsilon_{1}=\tau^{-1}.$ Since $r\leq 1$ again by Remark \ref{remsc}, we have 
\[\mathcal{L}_{\kappa}^{-}\big(D^{2}v,Dv\big)\leq 1~\text{in}~B_{2}.\]
Thus by Theorem \ref{lepsilon} we find  $v>1$ in $B_{1}$ as long as $\tau r^{2-\alpha}k\leq \epsilon_{0}.$ Which is the case provided we choose $\tilde{\kappa}_{0}=\kappa_{0}\tau^{-1}=\kappa_{0}\epsilon_{1},$ where $\kappa_{0}$ is from Lemma \ref{lepaa}. The required result follows once we scale back.
\end{proof}

\end{document}
